\section{Criterios de calidad y contenido de un requisito}
Un requisito de calidad debe aportar valor, estar ligado a un stakeholder concreto y situarse en un contexto claro. El handbook subraya que escribir requisitos no es un fin en s\'i mismo: su valor depende del beneficio que aportan frente al coste de especificarlos, validarlos, documentarlos y gestionarlos \cite[pp.~16--17]{cpre}.

Adem\'as, un requisito debe considerar las perspectivas de los stakeholders, porque RE trata de satisfacer sus deseos y necesidades, no solo las del usuario final. Eso obliga a identificar roles, priorizar su influencia y resolver conflictos entre puntos de vista distintos \cite[pp.~18--19]{cpre}.

Tambi\'en debe quedar claro el contexto del sistema, su l\'imite y su alcance. El handbook explica que los requisitos nunca aparecen aislados, sino dentro de un entorno relevante para entender el sistema; por eso hay que delimitar qu\'e est\'a dentro y fuera de alcance y qu\'e supuestos del contexto deben cumplirse \cite[pp.~22--24]{cpre}.

Por \'ultimo, un requisito debe ser expl\'icito, verificable y validado. El documento insiste en que la validaci\'on comprueba que las necesidades quedan cubiertas y que existe acuerdo entre las partes interesadas. En la gesti\'on posterior, adem\'as, conviene registrar identificaci\'on, prioridad, dependencias, origen, racionalidad y tipo del requisito \cite[pp.~109--111, 136--137]{cpre}.